% Figure: two point sources at the Rayleigh resolution limit. Each is an Airy
% profile (approximated here by a sinc^2-like / gaussian-shaped lobe); the
% peak of one falls on the first minimum of the other. The sum dips ~26%
% between the peaks, the classical Rayleigh "just resolved" criterion.
\begin{figure}[htbp]
\centering
\begin{tikzpicture}
\begin{axis}[
  width=12cm, height=6.6cm,
  xlabel={angular position (units of $1.22\lambda/D$)},
  ylabel={relative intensity},
  xmin=-2.6, xmax=2.6, ymin=0, ymax=1.15,
  axis lines=middle,
  xtick={-2,-1,0,1,2},
  ytick={0,0.5,1.0},
  tick label style={font=\scriptsize},
  label style={font=\small},
  samples=400, domain=-2.6:2.6,
  legend style={font=\scriptsize, at={(0.98,0.98)}, anchor=north east, draw=none},
]
% Airy-like lobe modeled as (2 J1(x)/x)^2 surrogate via (sin/.)^2 scaled.
% Place sources at -0.5 and +0.5 so peaks sit on each other's first zero (zero
% of the surrogate at distance 1.0). Use (sin(pi u)/(pi u))^2 with u offset.
\addplot[engblue, thick, dashed] {(sin(deg(pi*(x+0.5)))/(pi*(x+0.5)+1e-9))^2};
\addplot[engorange, thick, dashed] {(sin(deg(pi*(x-0.5)))/(pi*(x-0.5)+1e-9))^2};
\addplot[engred, very thick]
  {(sin(deg(pi*(x+0.5)))/(pi*(x+0.5)+1e-9))^2 + (sin(deg(pi*(x-0.5)))/(pi*(x-0.5)+1e-9))^2};
\legend{source A, source B, sum}
\end{axis}
\end{tikzpicture}
\caption[Rayleigh resolution criterion]{Two equal point sources at the
Rayleigh limit: the peak of each diffraction pattern falls on the first dark
ring of the other, a separation of $\theta_{\min} = 1.22\lambda/D$ for a
circular aperture of diameter $D$. The summed profile (solid) shows a
central dip of about a quarter below the peaks, the conventional threshold at
which two sources are called ``just resolved.'' Closer than this, the dip
fills in and the pair reads as a single blur. The lobes are drawn with a
$\sin^2$ surrogate for the Airy profile.}
\label{fig:vol04ch11:rayleigh-resolution}
\end{figure}
